% !TEX root = ../main.tex
\chapter{Phương pháp đề xuất}
\label{Chapter3}

Chương này trình bày kiến trúc CoQLLM, cơ sở thiết kế và quy trình huấn luyện của mô hình.
Xuất phát từ hạn chế của phép ánh xạ pointwise bằng MLP trong CoLLM, luận văn khảo sát
Q-Former như một cầu nối căn chỉnh có khả năng tổng hợp nhiều nguồn tín hiệu cộng tác trước
khi chuyển chúng thành soft token cho LLM. Các thành phần của mô hình được mô tả lần lượt từ
kiến trúc tổng thể, mô-đun MF, cầu nối Q-Former và cơ chế soft token đến các kỹ thuật huấn
luyện bổ sung và quy trình tối ưu ba giai đoạn.

\section{Kiến trúc tổng thể}
\label{sec:kien-truc-tong-quan}

\textbf{Bài toán.} Luận văn xem nhiệm vụ gợi ý dưới dạng dự đoán nhị phân. Với người dùng
$u$, lịch sử tương tác của người dùng và một sản phẩm mục tiêu $i$, mô hình ước lượng khả
năng người dùng thích sản phẩm đó và biểu diễn đầu ra của LLM dưới dạng câu trả lời ``Yes''
hoặc ``No''. Đây là bài toán phân loại nhị phân, được đánh giá chủ yếu bằng AUC và uAUC
(chi tiết ở Chương~\ref{Chapter4}).

CoQLLM gồm ba khối chính, minh họa ở Hình~\ref{fig:kien-truc}:

\begin{enumerate}
    \item \textbf{Mô-đun phân rã ma trận (MF)} cung cấp embedding cộng tác cho người dùng,
    sản phẩm mục tiêu và các sản phẩm trong lịch sử tương tác.

    \item \textbf{Cầu nối Q-Former} biến đổi mỗi véc-tơ cộng tác thành các soft token bằng một
    tập truy vấn học được (self-attention giữa các truy vấn và cross-attention tới véc-tơ cộng
    tác).

    \item \textbf{Mô hình ngôn ngữ lớn (LLM)} nhận các soft token cộng tác cùng prompt văn
    bản, kết hợp chúng qua self-attention và thực hiện dự đoán ``Yes/No''.
\end{enumerate}

\begin{figure}[htp]
\centering
\includegraphics[width=12cm]{images/coqllm-architecture.drawio.png}
\caption{Kiến trúc tổng quan của CoQLLM. Mô-đun MF cung cấp biểu diễn cộng tác; cầu nối
Q-Former căn chỉnh tín hiệu này thành các soft token được chèn vào prompt; LLM thực hiện dự
đoán. Chỉ Q-Former, các lớp chiếu và bộ chuyển thể LoRA được cập nhật theo từng giai đoạn
huấn luyện.}
\label{fig:kien-truc}
\end{figure}

Phần lớn tham số của hệ thống được giữ cố định trong quá trình huấn luyện. Cả MF lẫn backbone
LLM đều được đóng băng, trong khi Q-Former, các lớp chiếu và LoRA~\cite{lora2021} được cập
nhật ở các giai đoạn tương ứng. Tín hiệu cộng tác chỉ đi vào LLM qua soft token trong không
gian đầu vào, không thông qua cơ chế sinh trọng số theo từng mẫu. Thiết kế này tách tương đối
rõ vai trò của mô hình cộng tác, cầu nối căn chỉnh và mô hình ngôn ngữ, qua đó thuận lợi cho
việc phân tích đóng góp của từng thành phần.

\section{Mô-đun phân rã ma trận}
\label{sec:mf}

Mô-đun MF đóng vai trò cung cấp biểu diễn cộng tác đầu vào cho CoQLLM. Từ ma trận tương tác
người dùng--sản phẩm, mô hình học các embedding $e_u, e_i \in \mathbb{R}^{d}$ sao cho mức độ
tương thích giữa hai véc-tơ phản ánh sở thích của người dùng, đo bằng tích vô hướng
$e_u^\top e_i$. MF được huấn luyện bằng mục tiêu xếp hạng cặp BPR~\cite{bpr2009} và được đóng
băng trước khi tích hợp với các thành phần còn lại.

Ba nhóm tín hiệu được trích xuất từ MF gồm: (i)~embedding người dùng $e_u$; (ii)~embedding
sản phẩm mục tiêu $e_i$; và (iii)~embedding của các sản phẩm trong lịch sử tương tác
$\{e_{i'} : i' \in \mathcal{H}_u\}$.

Việc huấn luyện MF tách biệt khỏi Q-Former và LLM nhằm giữ cố định nguồn tín hiệu cộng tác
trong các thí nghiệm phía sau. Cách thiết kế này giúp giảm sự đan xen giữa thay đổi của mô
hình cộng tác và thay đổi của cầu nối căn chỉnh, từ đó thuận lợi hơn khi phân tích nguyên
nhân của các chênh lệch thực nghiệm. Trong phạm vi luận văn, MF được huấn luyện trên phần dữ
liệu huấn luyện tương ứng của từng tập dữ liệu và sau đó giữ nguyên trong các giai đoạn tích
hợp.

Việc giữ MF đóng băng trong các giai đoạn sau nhất quán với cấu hình mặc định của
CoLLM~\cite{collm2023}. Trong khảo sát của CoLLM, việc tiếp tục tinh chỉnh mô hình cộng tác
cùng với lớp chiếu có thể cải thiện nhẹ kết quả trong phần lớn trường hợp, nhưng làm tăng
đáng kể chi phí huấn luyện (khoảng năm lần trên MovieLens-1M) và làm chậm hội tụ; vì vậy đóng
băng là một đánh đổi hợp lý giữa hiệu năng và chi phí. Bên cạnh đó, một mô hình cộng tác đã
được huấn luyện tốt còn đóng vai trò tiên nghiệm (prior) và ràng buộc cho quá trình học soft
token. Quan trọng hơn đối với mục tiêu phân tích của luận văn, việc cố định MF cho phép quy
phần trần của uAUC về chất lượng embedding của MF ở nhóm người dùng/sản phẩm nguội
(Mục~\ref{ssec:user-group}), thay vì để đóng góp của MF và của cầu nối căn chỉnh trộn lẫn với
nhau. Việc mở băng và tinh chỉnh đồng thời MF được ghi nhận như một hướng khảo sát trong
tương lai.

\section{Cầu nối Q-Former}
\label{sec:qformer}

\subsection{Động cơ lựa chọn Q-Former}
\label{ssec:ly-do-qformer}

Mô-đun ánh xạ trong CoLLM~\cite{collm2023} có dạng MLP hai lớp:
\begin{equation*}
g_\phi(\mathbf{x}) = \mathbf{W}_2\,\mathrm{ReLU}(\mathbf{W}_1\mathbf{x} + \mathbf{b}_1)
+ \mathbf{b}_2,
\end{equation*}
với $g_\phi: \mathbb{R}^d \rightarrow \mathbb{R}^{d_{\text{LLM}}}$. Phép biến đổi này được áp
dụng độc lập trên từng embedding cộng tác và mỗi embedding tạo ra một soft token. Cần lưu ý
rằng MLP đã là một phép biến đổi phi tuyến; hạn chế mà luận văn quan tâm không nằm ở tính
tuyến tính hay phi tuyến, mà ở cách thông tin được tổng hợp.

Thứ nhất, phép ánh xạ có tính pointwise: đầu ra ứng với một embedding được tính độc lập với
các embedding còn lại trong cùng mẫu. Vì vậy, embedding người dùng, sản phẩm mục tiêu và lịch
sử không được phối hợp ở phía mô-đun ánh xạ trước khi đi vào LLM. Thứ hai, mỗi embedding được
cô đọng thành một soft token duy nhất, nên mô-đun không có một cơ chế riêng để học cách phân
bổ mức độ chú ý khác nhau cho các nguồn tín hiệu cộng tác theo từng mẫu.

Q-Former cung cấp một cơ chế khác. Thay vì một phép chiếu cố định cho ra đúng một token, mỗi
véc-tơ cộng tác được một \textit{tập truy vấn học được} khai triển thành nhiều biểu diễn: các
truy vấn có khởi tạo riêng và phối hợp với nhau qua self-attention, nên có thể chuyên biệt hóa
để nắm các khía cạnh khác nhau của cùng một tín hiệu cộng tác. Đây là điểm khác biệt cốt lõi
so với một hàm pointwise duy nhất.

Cần nói rõ giới hạn của lập luận này để tránh phóng đại. Lợi thế của Q-Former ở đây không đến
từ khả năng nén như trong thị giác: mỗi sản phẩm được mã hóa từ \textit{một} véc-tơ cộng tác
duy nhất, nên cross-attention chủ yếu đóng vai một phép đọc có học đối với véc-tơ đó (chi tiết
ở Mục~\ref{ssec:qformer-kien-truc}). Sức mạnh của cầu nối vì thế nằm ở tập truy vấn học được,
sự phối hợp qua self-attention và độ sâu của phép biến đổi, chứ không nằm ở việc gộp nhiều
nguồn tín hiệu bên trong Q-Former.

\subsection{Q-Former cho bài toán gợi ý}
\label{ssec:qformer-kien-truc}

CoQLLM xây dựng cầu nối dựa trên kiến trúc Q-Former của BLIP-2~\cite{blip2_2023}. Mô hình sử
dụng $Q = 8$ véc-tơ truy vấn học được $\mathbf{Q} \in \mathbb{R}^{Q \times d_{\text{model}}}$,
với $d_{\text{model}} = 768$. Nhánh văn bản được khởi tạo từ \texttt{bert-base-uncased} và
được sử dụng trong các mục tiêu căn chỉnh sản phẩm--văn bản ở giai đoạn tiền huấn luyện biểu
diễn (Mục~\ref{ssec:stage1}).

Trong mỗi lớp Q-Former, các truy vấn đi qua hai cơ chế attention:

\begin{enumerate}
    \item \textit{Self-attention} giữa các truy vấn với nhau, cho phép các truy vấn trao đổi
    thông tin và phối hợp vai trò biểu diễn, mỗi truy vấn đảm nhận một khía cạnh khác nhau của
    tín hiệu cộng tác.

    \item \textit{Cross-attention} từ các truy vấn đến biểu diễn cộng tác của sản phẩm đang được
    mã hóa: các truy vấn cùng đọc một véc-tơ cộng tác đầu vào và biến đổi nó thành nhiều kênh
    biểu diễn khác nhau.
\end{enumerate}

Đầu ra là khối truy vấn $\mathbf{Z} = \text{cf\_q} \in \mathbb{R}^{B \times Q \times
d_{\text{model}}}$, trong đó $B$ là kích thước batch và $Q = 8$.

Trong CoQLLM, mỗi sản phẩm (sản phẩm mục tiêu và từng sản phẩm trong lịch sử) được Q-Former mã
hóa \textit{độc lập} từ một véc-tơ cộng tác thành $Q$ soft token. Tín hiệu người dùng không đi
vào Q-Former dưới dạng đầu vào cross-attention mà được đưa vào qua điều kiện hóa truy vấn
(Mục~\ref{ssec:user-conditioned}); còn việc kết hợp thông tin \textit{giữa} các sản phẩm được
thực hiện ở lớp self-attention của LLM trên toàn bộ soft token, chứ không phải bên trong cầu
nối.

Điều này dẫn tới một khác biệt so với BLIP-2. Trong thị giác, cross-attention nén hàng trăm
patch thành số truy vấn nhỏ hơn nhiều, nên đóng vai trò nén chính. Trong CoQLLM, do đầu vào chỉ
là một véc-tơ cộng tác cho mỗi lần mã hóa, cross-attention gần như một phép đọc có học đối với
véc-tơ đó; điều tạo ra $Q$ biểu diễn khác nhau chính là tập truy vấn học được (mỗi truy vấn
khởi tạo riêng) cùng self-attention phối hợp giữa chúng. Vì vậy, luận văn xem \textit{tập truy
vấn học được và self-attention} -- chứ không phải khả năng nén hay việc gộp nhiều nguồn bên
trong Q-Former -- là thành phần quan trọng của cầu nối căn chỉnh.

\subsection{Tích hợp biểu diễn cộng tác qua soft token}
\label{ssec:hai-duong}

Đầu ra $\mathbf{Z}$ của Q-Former được chiếu từ không gian $d_{\text{model}}$ sang không gian
embedding của LLM bằng lớp tuyến tính $W_{\text{proj}} \in \mathbb{R}^{d_{\text{model}} \times
d_{\text{LLM}}}$. Kết quả là $Q$ soft token được chèn vào prompt tại vị trí dành cho biểu diễn
cộng tác.

Phần văn bản của prompt đánh dấu lịch sử và sản phẩm mục tiêu bằng các token chuyên dụng. Một
token định danh riêng được sử dụng cho vị trí soft token nhằm tránh phụ thuộc vào
\verb|<unk>| hoặc các token có chức năng khác trong từ vựng gốc. Sau lớp chiếu, một lớp chuẩn
hóa (LayerNorm) được áp dụng để giữ thang độ của soft token tương thích hơn với embedding văn
bản của LLM.

Nhờ cơ chế này, tín hiệu cộng tác được đưa vào cùng không gian đầu vào với văn bản mà không
cần thay đổi trực tiếp trọng số backbone LLM theo từng mẫu. Đây là điểm khác với
CoRA~\cite{cora2024}, nơi biểu diễn cộng tác được sử dụng để sinh cập nhật trọng số hạng thấp
cho các lớp attention. CoQLLM chủ động giới hạn đường tích hợp ở soft token để giữ LLM hoàn
toàn đóng băng và tập trung phân tích vai trò của cầu nối căn chỉnh. Trên nền đường soft token
này, luận văn bổ sung hai kỹ thuật huấn luyện -- điều kiện hóa truy vấn theo người dùng
(Mục~\ref{ssec:user-conditioned}) và mục tiêu bảo toàn thứ tự trong từng người dùng
(Mục~\ref{ssec:align-rank}) -- nhằm cải thiện lần lượt AUC và uAUC.

\subsection{Điều kiện hóa truy vấn theo người dùng}
\label{ssec:user-conditioned}

Trong cấu hình cơ sở, tập truy vấn $\mathbf{Q}$ là tham số dùng chung cho mọi mẫu. Để tăng khả
năng thích nghi theo người dùng, luận văn bổ sung một phép dịch chuyển phụ thuộc vào embedding
người dùng trước khi đưa truy vấn vào Q-Former:

\begin{equation} \label{eq:user-cond}
\tilde{\mathbf{Q}} = \mathbf{Q} + g_u\big( e_u \big),
\end{equation}

trong đó $g_u$ là một lớp chiếu tuyến tính từ không gian MF sang không gian truy vấn, được
\textit{khởi tạo bằng không} (zero-init). Với cách khởi tạo này, ở bước đầu cơ chế không làm
thay đổi truy vấn gốc, tức $\tilde{\mathbf{Q}} = \mathbf{Q}$; mức dịch chuyển chỉ được học khi
mang lại lợi ích cho hàm mục tiêu trong quá trình huấn luyện.

Về mặt trực giác, thành phần $g_u(e_u)$ tạo ra một mức điều chỉnh riêng theo người dùng.
Vì phần dịch chuyển \textit{thay đổi giữa các người dùng} nhưng \textit{không đổi trong phạm
vi một người dùng}, cơ chế này có thể nâng khả năng phân biệt giữa các người dùng và do đó tác
động tích cực đến AUC toàn cục, trong khi tác động lên uAUC -- vốn bất biến với dịch chuyển mức
nền riêng theo từng người dùng -- có thể nhỏ hơn. Giả thuyết này được đối chiếu với kết quả
thực nghiệm ở Chương~\ref{Chapter4}. Vì $g_u$ chỉ xuất hiện khi bật cờ và được khởi tạo bằng
không, cơ chế tương thích với checkpoint cũ và chỉ cần huấn luyện lại Bước 2 của Giai đoạn 3.

\subsection{Hàm mất mát bảo toàn thứ hạng theo người dùng}
\label{ssec:align-rank}

Hàm mất mát mặc định trong giai đoạn tinh chỉnh CTR là entropy chéo nhị phân (BCE) trên nhãn
``Yes''/``No''. BCE tối ưu khả năng phân loại từng mẫu nhưng không trực tiếp tối ưu thứ tự
tương đối giữa các sản phẩm của cùng một người dùng, trong khi uAUC đo chính đặc tính xếp hạng
nội bộ này.

Để bổ sung tín hiệu xếp hạng, luận văn sử dụng một mục tiêu BPR theo cặp trong phạm vi từng
người dùng (\textit{align-rank loss}), áp trực tiếp lên soft token cộng tác đã căn chỉnh
(trước khi vào LLM) thông qua một đầu ra phụ cũng khởi tạo bằng không:

\begin{equation} \label{eq:align-rank}
\mathcal{L} = \mathcal{L}_{\text{BCE}}
  + \lambda \sum_{u} \sum_{(i^{+},\, i^{-}) \in \mathcal{P}_u}
    -\log \sigma\!\Big( \tfrac{1}{\tau}\big( s_u(i^{+}) - s_u(i^{-}) \big) \Big),
\end{equation}

trong đó $s_u(\cdot)$ là điểm xếp hạng được sinh từ soft token cộng tác đã căn chỉnh,
$\mathcal{P}_u$ là tập các cặp dương--âm của cùng người dùng $u$, $\sigma$ là hàm sigmoid,
$\tau$ là nhiệt độ và $\lambda$ là trọng số của thành phần xếp hạng.

Mục tiêu phụ được áp trực tiếp lên biểu diễn cộng tác sau căn chỉnh, thay vì chỉ lên logit
``Yes/No'' của LLM, nhằm buộc bản thân cầu nối giữ đúng thứ tự giữa sản phẩm dương và âm trong
từng người dùng. Để tạo đủ cặp dương--âm trong một batch, quá trình huấn luyện sử dụng lấy mẫu
theo nhóm người dùng (\textit{user-grouped batch}) và mục tiêu được kích hoạt ở Bước 2 của Giai
đoạn 3 (nơi Q-Former và lớp chiếu được huấn luyện). Do đầu ra phụ khởi tạo bằng không, các
epoch đầu gần như trung tính và tác động của mục tiêu xếp hạng hiện dần về sau. Đây là một siêu
tham số phụ thuộc tập dữ liệu: nó cải thiện uAUC trên tập còn nhiều dư địa xếp hạng trong từng
người dùng (Amazon-Book) và gần như trung tính khi uAUC đã bão hòa (MovieLens-1M).

\section{Biểu diễn cộng tác đa token}
\label{sec:lua-chon-thiet-ke}

Do Q-Former sử dụng $Q = 8$ truy vấn học được, đầu ra của nó là $Q$ véc-tơ. CoQLLM chiếu toàn
bộ $Q$ véc-tơ này qua lớp chiếu và đưa vào LLM như $Q$ soft token riêng biệt cho mỗi sản phẩm
(sản phẩm mục tiêu và từng sản phẩm trong lịch sử). Đây là \textit{hệ quả trực tiếp} của việc
dùng nhiều truy vấn, không phải một mô-đun được thiết kế thêm: giữ nguyên cả $Q$ véc-tơ cho
phép mỗi truy vấn truyền một khía cạnh biểu diễn khác nhau tới LLM, thay vì phải hợp nhất về
một véc-tơ duy nhất trước khi đi vào mô hình ngôn ngữ.

Điểm này tạo một khác biệt so với CoLLM~\cite{collm2023}, vốn cô đọng mỗi véc-tơ cộng tác
thành đúng một soft token. Với cùng số nguồn cộng tác, CoQLLM đưa vào LLM số soft token gấp $Q$
lần, tức băng thông biểu diễn cao hơn. Luận văn ghi nhận đây là một \textit{yếu tố gây nhiễu}
(confound) cần lưu ý khi so sánh với CoLLM, chứ không xem việc dùng nhiều token là một đóng góp
riêng.

Việc tách riêng ảnh hưởng của số lượng soft token (một token so với $Q$ token) dưới cùng một
backbone chưa được thực hiện trong luận văn; đây là một ablation còn để ngỏ, được ghi nhận minh
bạch ở Mục~\ref{sec:han-che}.

\section{Huấn luyện ba giai đoạn}
\label{sec:ba-giai-doan}

CoQLLM được huấn luyện theo ba giai đoạn nối tiếp, theo quy trình của CoLLM~\cite{collm2023}
với điều chỉnh cho phù hợp với kiến trúc Q-Former. Quy trình kế thừa tinh thần tách biệt giữa
học biểu diễn cộng tác, căn chỉnh với không gian ngôn ngữ và tinh chỉnh cho nhiệm vụ gợi ý.

\subsection{Tiền huấn luyện biểu diễn}
\label{ssec:stage1}

Giai đoạn đầu tập trung giúp Q-Former học mối quan hệ giữa biểu diễn cộng tác và thông tin văn
bản, thông qua bốn mục tiêu:

\begin{itemize}
    \item \textit{Căn chỉnh sản phẩm--văn bản (ITC)}: tăng độ tương đồng giữa biểu diễn
    Q-Former của một sản phẩm (từ embedding MF) và biểu diễn văn bản tương ứng (tên phim, thể
    loại). Đây là mục tiêu căn chỉnh chính.

    \item \textit{So khớp sản phẩm--văn bản (ITM)}: phân loại nhị phân một cặp sản phẩm--văn
    bản là tương ứng hay không tương ứng.

    \item \textit{Sinh văn bản có điều kiện (ITG)}: sử dụng biểu diễn sản phẩm để hỗ trợ sinh
    mô tả văn bản.

    \item \textit{Tương phản người dùng--sản phẩm}: đưa các cặp tương tác dương lại gần nhau và
    tách các cặp âm trong không gian biểu diễn.
\end{itemize}

Một vấn đề thực tế là sự mất cân bằng giữa số cặp người dùng--sản phẩm và số cặp sản phẩm--văn
bản. Trên MovieLens-1M, số cặp người dùng--sản phẩm dương khoảng $888$ nghìn, trong khi số cặp
sản phẩm--văn bản chỉ khoảng $3$ nghìn (tỉ lệ khoảng $0{,}34\%$). Nếu lấy mẫu trực tiếp từ toàn
bộ dữ liệu, các mục tiêu ITC/ITM/ITG có thể được kích hoạt quá ít trong từng batch -- điều kiện
``ít nhất hai mẫu cùng loại'' hiếm khi được thỏa -- và gradient từ mục tiêu người dùng--sản
phẩm chiếm ưu thế.

Để giảm mất cân bằng, luận văn áp dụng hai điều chỉnh, tổng hợp ở
Bảng~\ref{tab:balance-stage1}:

\begin{table}[ht]
\caption{Hai điều chỉnh cân bằng dữ liệu Giai đoạn 1}
\label{tab:balance-stage1}
\begin{center}
\begin{tabular}{|L{4.5cm}|L{2.5cm}|L{5cm}|}
\hline
\textbf{Điều chỉnh} & \textbf{Mức áp dụng} & \textbf{Mục đích} \\
\hline
Giới hạn số cặp người dùng--sản phẩm ở mức $20{.}000$ & Dữ liệu & Tăng tỉ lệ mẫu sản
phẩm--văn bản trong batch, giúp các mục tiêu ITC/ITM/ITG được kích hoạt thường xuyên hơn \\
\hline
Giảm trọng số mất mát người dùng--sản phẩm từ $1{,}0$ xuống $0{,}5$ & Hàm mất mát & Cân bằng
cường độ gradient giữa các mục tiêu \\
\hline
\end{tabular}
\end{center}
\end{table}

Danh sách cặp người dùng--sản phẩm kế thừa thứ tự sắp theo thời gian của dữ liệu gốc, nên việc
lấy phần đầu danh sách sẽ thiên lệch về các tương tác sớm. Để tránh thiên lệch này, danh sách
được xáo trộn ngẫu nhiên với hạt giống cố định trước khi cắt, biến giới hạn $20{.}000$ thành
một mẫu ngẫu nhiên đồng nhất mà vẫn bảo đảm tính tái lập.

\subsection{Tiền huấn luyện sinh}
\label{ssec:stage2}

Ở giai đoạn thứ hai, Q-Former và lớp chiếu được huấn luyện với mục tiêu sinh ngôn ngữ. Từ biểu
diễn cộng tác đã qua Q-Former, mô hình học tạo ra các soft token mà LLM đóng băng có thể khai
thác để sinh văn bản mô tả sản phẩm. Có thể xem đây là bước căn chỉnh đầu ra của Q-Former với
không gian embedding mà LLM sử dụng trong quá trình sinh.

\subsection{Tinh chỉnh CTR}
\label{ssec:stage3}

Giai đoạn cuối theo quy trình hai bước của CoLLM~\cite{collm2023}.

\textbf{Bước 1 -- Tinh chỉnh LoRA chỉ với thông tin văn bản.} LoRA được huấn luyện bằng prompt
chỉ chứa tên sản phẩm, thể loại và lịch sử ở dạng văn bản, chưa đưa tín hiệu cộng tác vào. Mục
tiêu của bước này là giúp LLM thích nghi với nhiệm vụ dự đoán nhị phân từ thông tin ngôn ngữ và
ổn định hóa mô hình trước khi tiếp nhận soft token cộng tác.

\textbf{Bước 2 -- Tinh chỉnh cầu nối căn chỉnh với tín hiệu cộng tác.} Mô hình được khởi tạo từ
checkpoint của Bước 1. LoRA được giữ cố định, trong khi Q-Former và các lớp chiếu được cập nhật
để học cách đưa biểu diễn cộng tác vào LLM dưới dạng soft token -- tách bạch phần ``ngôn ngữ''
(LoRA, học ở Bước 1) khỏi phần ``cộng tác'' (cầu nối căn chỉnh, học ở Bước 2). Đây cũng là bước
duy nhất kích hoạt hai kỹ thuật điều kiện hóa truy vấn theo người dùng
(Mục~\ref{ssec:user-conditioned}) và, khi cấu hình cho phép, mục tiêu bảo toàn thứ hạng trong
từng người dùng (Mục~\ref{ssec:align-rank}). Nhãn giám sát là ``Yes''/``No'' tương ứng với nhãn
CTR, và checkpoint tốt nhất được chọn theo uAUC trên tập kiểm định.

So với mô-đun chiếu MLP của CoLLM, Q-Former làm tăng chi phí huấn luyện do có nhiều lớp
attention và tham số hơn. Vì tín hiệu cộng tác chỉ đi qua đường soft token và không sinh trọng
số theo từng mẫu, LLM được giữ nguyên hoàn toàn khi suy luận. Luận văn xem đây là một đánh đổi
có chủ đích để khảo sát khả năng căn chỉnh biểu diễn tốt hơn, thay vì giả định chi phí tăng
thêm là ``không đáng kể''. Chi tiết về số tham số và thời gian huấn luyện được trình bày ở
Chương~\ref{Chapter4}.

\section{Tổng kết chương}
\label{sec:ket-chuong-3}

Chương này đã trình bày CoQLLM như một kiến trúc tích hợp tín hiệu cộng tác vào LLM thông qua
cầu nối Q-Former. Khác với phép ánh xạ pointwise bằng MLP của CoLLM, Q-Former tổng hợp nhiều
nguồn tín hiệu bằng các truy vấn học được và attention trước khi chuyển chúng thành soft token.
Trên nền kiến trúc này, luận văn bổ sung điều kiện hóa truy vấn theo người dùng và mục tiêu bảo
toàn thứ hạng trong từng người dùng -- cải thiện lần lượt AUC và uAUC. Chương cũng ghi nhận
biểu diễn đa token như một hệ quả trực tiếp của thiết kế nhiều truy vấn, đồng thời là yếu tố
cần lưu ý khi so sánh với CoLLM. Chương tiếp theo đánh giá các thiết kế trên thông qua kết quả
thực nghiệm.
